\section{What spectral clustering implies for linear solves}
\label{sec:clustered-solves}

Large spectral condition numbers need not force the usual
$O(\sqrt\kappa)$ conjugate-gradient iteration bound to be sharp.
Classical cluster-sensitive estimates exploit the spectrum more
precisely; see \citet{axelsson1994}. We give the elementary polynomial
construction needed for \cref{thm:lp-all-spectra}. This is an
exact-arithmetic statement about energy-norm error, not a
finite-precision residual theorem.

\begin{proposition}[An explicit two-interval CG bound]
\label{prop:clustered-cg}
Suppose $H\succ0$ and
\[
 \operatorname{spec}(H)\subseteq[a_1,b_1]\cup[a_2,b_2],\qquad
 0<a_1\leq b_1<a_2\leq b_2.
\]
Put $\rho_i=b_i/a_i$ and $\Lambda=b_2/a_1$. For
$0<\delta\leq1/2$ there is a real polynomial $q$, with $q(0)=1$,
such that $\max_{\operatorname{spec}(H)}|q|\leq\delta$ and
\begin{equation}
 \deg q=O\left(
 \sqrt{\rho_1\rho_2}\log(4\Lambda)\log(1/\delta)
 +\sqrt{\rho_2}\log(1/\delta)\right).
 \label{eq:clustered-cg-degree}
\end{equation}
Exact-arithmetic CG for $Hu=r$, from any initial vector $u_0$,
therefore achieves
$\norm{u-u_k}_H\leq\delta\norm{u-u_0}_H$ within this many
iterations, or within $\dim H$ iterations if that is smaller.
\end{proposition}

\begin{proof}
For $0<a<b$ the normalized Chebyshev polynomial
\[
 r_m(x;a,b)=
 \frac{T_m((a+b-2x)/(b-a))}{T_m((a+b)/(b-a))}
\]
has value one at zero and satisfies
\begin{equation}
 \max_{[a,b]}|r_m|\leq
 2\left(\frac{\sqrt{b/a}-1}{\sqrt{b/a}+1}\right)^m
 \leq2\exp\left(-\frac{2m}{\sqrt{b/a}}\right).
 \label{eq:chebyshev-bound}
\end{equation}
The first bound follows from $|T_m|\leq1$ on $[-1,1]$ and
$T_m(t)=(w^m+w^{-m})/2$ for
$t=(w+w^{-1})/2>1$. The second uses
$\log((s+1)/(s-1))\geq2/s$ for $s>1$.

Choose $p_2=r_{m_2}(\cdot;a_2,b_2)$ with
$m_2=\lceil\sqrt{\rho_2}\rceil$. Its magnitude is at most $1/2$
on the upper interval, and at most one on $[0,a_2]$, because
$T_{m_2}(t)$ is positive and increasing for $t\geq1$. If
$a_2=b_2$, use $p_2(x)=1-x/a_2$, $m_2=1$ instead.
Choose $r_1=r_{m_1}(\cdot;a_1,b_1)$ with
$m_1=\lceil\tfrac12\sqrt{\rho_1}\log(2/\delta)\rceil$,
so that $|r_1|\leq\delta$ on the lower interval. If $a_1=b_1$,
use $r_1(x)=1-x/a_1$ and $m_1=1$.

All $m_1$ roots $\theta_j$ of $r_1$ lie in $[a_1,b_1]$, and
$r_1(0)=1$, so its product representation gives
\[
 |r_1(x)|=\prod_{j=1}^{m_1}|1-x/\theta_j|
 \leq\Lambda^{m_1}\qquad(a_2\leq x\leq b_2).
\]
Let
\[
 j=\left\lceil m_1\log_2(4\Lambda)+\log_2(1/\delta)\right\rceil,
 \qquad q=r_1p_2^j.
\]
On the lower interval $|q|\leq\delta$; on the upper interval
$|q|\leq\Lambda^{m_1}2^{-j}\leq\delta$. Its degree
$m_1+jm_2$ gives \eqref{eq:clustered-cg-degree}.
Finally, exact-arithmetic CG minimizes the energy-norm error over
vectors $u-u_k=q_k(H)(u-u_0)$ with $q_k(0)=1$ and degree at most
$k$. The spectral theorem bounds the relative energy error of our
candidate by $\max_{\operatorname{spec}(H)}|q|$. The usual exact
termination in dimension $\dim H$ gives the stated alternative.
\end{proof}

For each fixed barrier in \cref{thm:lp-all-spectra}, the two interval
ratios are bounded and $\Lambda=\Theta(g^{-2})$ when $f>0$.
Thus \cref{prop:clustered-cg} gives
$O(\log(1/g)\log(1/\delta))$ iterations on a sufficiently small
tail. If $f=0$, there is one interval with bounded condition number,
and the ordinary Chebyshev estimate gives $O(\log(1/\delta))$.
These are matrix--vector-product counts. They do not include the cost
of producing those products, and floating-point loss of conjugacy can
affect both attainable error and iteration counts. A recursively
updated residual should therefore be checked against a freshly
computed $r-Hu_k$ in any numerical use of the bound.

\subsection{Quantum access and output are separate assumptions}

A quantum linear-system algorithm commonly receives a unitary whose
specified block equals $H/\alpha$, up to a controlled encoding error,
where $\alpha\geq\norm H$ is its normalization. This is called a
\emph{block encoding}. The inverse scale seen by that encoding is
\begin{equation}
 \kappa_{\mathrm{enc}}=\alpha\norm{H^{-1}},
 \label{eq:encoding-condition}
\end{equation}
which can exceed $\kappa(H)$. Such an algorithm also needs a procedure
preparing the normalized actual right-hand side and a specified output
task, often preparation of the normalized solution rather than all
its entries. These assumptions are part of the problem, not
consequences of Hessian geometry.

Quantum singular value transformation implements suitably bounded
polynomials of a normalized encoded matrix
\citep{gilyen2019}. The polynomial must meet the boundedness contract
on the whole interval $[-1,1]$, including points absent from the
spectrum. In contrast, \cref{prop:clustered-cg} controls its polynomial
only on the two spectral intervals; its values in their intervening
gap are unrestricted. That construction alone therefore supplies no
quantum implementation with its CG degree as the query cost.

The distinction is consistent with the established lower bounds and
structured positive results of \citet{orsucci2021}. Their
Proposition 6 proves a worst-case $\Omega(\min\{\kappa,N\})$
query lower bound for positive-definite quantum linear systems of
dimension $N$, with the specified matrix and state-preparation
oracles. The linear-in-$\kappa$ regime consequently requires
$N$ at least of that order. Such a worst-case family is not by itself
an accuracy lower bound along one fixed-dimensional central path.
Their accelerated positive-definite constructions impose additional
normalization, decomposition, or state-overlap conditions; merely
writing a Gram factorization of $H$ does not establish a generic
square-root dependence on $\kappa$.

There is also an important right-hand-side restriction. The geometric
theorems here concern all eigenvalues, whereas exact path-following
forces the particular vector in \eqref{eq:exact-newton-rhs}.
The window residual guarantee \eqref{eq:conditional-window-residual}
can exploit that vector only under its stated projection and solve
assumptions. Moreover, \cref{prop:oscillatory-barrier} shows that a
small discarded right-hand-side fraction can carry almost all of the
Euclidean solution norm. Thus a residual guarantee cannot silently be
substituted for a quantum solution-state guarantee.

Finally, Hessian conditioning is not an obstruction to every quantum
approach to optimization. \citet{apersgribling2026} give quantum
speedups for tall LPs with row access, using spectral approximations
and gradient estimation to return an explicit feasible approximate
solution. Their method avoids the Hessian-condition-number dependence
of a direct quantum linear-system substitution. The results here
describe a fixed metric and its spectrum; they neither rule out that
algorithm nor prove an end-to-end classical--quantum runtime
separation.
